\section{Methodik}
\label{sec:methodik}

Gegenstand der vorliegenden Untersuchung ist die Evaluation eines sokratischen
KI\-Tutors, der auf einer Retrieval\-Augmented\-Generation\-Pipeline (RAG) aufsetzt.
Ziel ist es, das Zusammenspiel von Prompting und Retrieval\-Konfiguration systematisch
zu bewerten. Die Untersuchung folgt einem explorativen, experimentellen
Evaluationsdesign, das zwei Forschungsfragen getrennt adressiert: den Einfluss
verschiedener sokratischer Prompts (RQ1) sowie den Einfluss verschiedener
RAG\-Konfigurationen (RQ2). Beide Fragen werden über dieselbe technische
Evaluationsarchitektur beantwortet, jedoch in zwei voneinander unabhängigen
Auswertungslinien, sodass jeweils nur ein Faktor variiert wird.

\subsection{Untersuchungsdesign}
\label{subsec:design}

Die Evaluation ist als automatisierte, mehrstufige Dialogbewertung angelegt. Für jedes
Testszenario wird ein Multi\-Turn\-Gespräch zwischen einer simulierten studierenden
Person und dem RAG\-gestützten Tutor erzeugt und anschließend anhand eines definierten
Metrik\-Sets bewertet. Die Wahl eines simulationsbasierten Multi\-Turn\-Designs ergibt
sich unmittelbar aus dem Untersuchungsgegenstand: Da sokratisches Tutoring ein
dialogisches, über mehrere Gesprächsrunden verlaufendes Vorgehen ist, lässt es sich
nicht anhand einzelner Frage\-Antwort\-Paare, sondern nur im Verlauf eines Gesprächs
angemessen beurteilen \cite{QUELLE_SOKRATIK}.

Zur klaren Trennung der Forschungsfragen wurden zwei separate Evaluationslinien
umgesetzt. In der ersten Linie (RQ1) wird bei konstant gehaltener RAG\-Konfiguration
der System\-Prompt variiert. Bewertet wird die \emph{\texttt{no\_Prompt}
\texttt{stock\_prompt}tät} des
Gesprächsverlaufs. In der zweiten Linie (RQ2) wird bei konstant gehaltenem Prompt die
RAG\-Konfiguration variiert und die \emph{Retrieval\-Qualität} bewertet. Durch die
Beschränkung auf jeweils einen variierten Faktor je Linie können die beobachteten
Effekte den Prompts bzw. den Konfigurationen zugeordnet werden.

\subsection{Untersuchungsgegenstand und Datengrundlage}
\label{subsec:gegenstand}

Das zu evaluierende System (System under Test) ist eine OpenWebUI\-kompatible
RAG\-Pipeline, die einen Qdrant\-Vektorspeicher, ein Cross\-Encoder\-Reranking sowie
Moodle\-Metadaten anbindet. Als generierendes Sprachmodell nutzt die Pipeline intern
das Modell \texttt{gemma\-4\-26B}. Der fachliche Wissenskorpus liegt in zwei
Qdrant\-Collections vor, die Vorlesungs\- bzw. Labormaterial des Fachgebiets
Elektronik der Hochschule enthalten.

Die Datengrundlage der Evaluation bilden nicht vorab gesammelte Dialoge, sondern
synthetisch erzeugte Gespräche. Grundlage jedes Gesprächs ist ein Testszenario,
das sich aus der Kombination von Themengebiet, Kompetenzniveau der studierenden Person
(Anfänger bzw. Fortgeschritten) und einem charakteristischen Studierenden\-Verhalten
(etwa kooperativ, die Lösung einfordernd oder inhaltlich selbstsicher\-fehlerhaft)
zusammensetzt. Insgesamt ergeben sich daraus 24 Szenarien je Auswertungslinie. Jedes Szenario wird als \emph{Golden} modelliert, das die
Ausgangssituation, eine Beschreibung der zu simulierenden Persona sowie eine feste
Initialfrage enthält. Die Initialfrage wird bewusst wortgleich vorgegeben und nicht
generiert, um über alle Prompts und Konfigurationen hinweg denselben Gesprächseinstieg
und damit die Vergleichbarkeit der Läufe sicherzustellen; die daran anschließenden
Folgefragen werden hingegen dynamisch simuliert.

% =========================================================
% TABELLE 1 · Szenario-Dimensionen
% =========================================================
\begin{table}[htbp]
  \centering
  \caption{Dimensionen der Testszenarien. Die Kombination aus Themengebiet,
  Kompetenzniveau und Studierenden-Verhalten ergibt die simulierten Gespr\"ache.}
  \label{tab:szenario-dimensionen}
  \begin{tabular}{ll}
    \toprule
    \textbf{Dimension} & \textbf{Auspr\"agungen} \\
    \midrule
    Themengebiet & Sperrspannung/Durchbruch der Z-Diode; \\
                 & Spannungsstabilisierung mit Z-Diode; \\
                 & Ohmsches Gesetz und Maschenregel \\
    \addlinespace
    Kompetenzniveau & Anf\"anger, Fortgeschritten \\
    \addlinespace
    Studierenden-Verhalten & kooperativ; gibt schnell auf; \\
                           & fordert die L\"osung ein; \\
                           & \"au{\ss}ert eine Vermutung; \\
                           & selbstsicher, aber inhaltlich falsch \\
    \bottomrule
  \end{tabular}
\end{table}
 
\subsection{Evaluationsarchitektur}
\label{subsec:architektur}

Da der Untersuchungsgegenstand ein Dialog ist, sind neben dem zu bewertenden System zwei
weitere Rollen erforderlich. Die Architektur umfasst somit drei klar getrennte Modelle.
Der Tutor ist die RAG\-Pipeline {gemma\-4\-26B} und beantwortet jeden
Gesprächszug. Der Studenten\-Simulator übernimmt die Rolle der studierenden
Person und erzeugt auf Basis von Szenario, Persona und der jeweils vorangegangenen
Tutor\-Antwort die Folgefragen. Ohne diese Rolle entstünde kein Dialog, sondern
lediglich ein einzelnes Frage\-Antwort\-Paar. Als Simulator wird das Modell
\texttt{gpt\-4o\-mini} eingesetzt, das auf dem verwendeten Endpunkt einen Alias für
\texttt{gemma\-4\-26B} darstellt und somit auf demselben Basismodell wie die
RAG\-Pipeline beruht. Der Judge bewertet den fertigen Dialog. Hierfür wird mit
gpt\-4.1 bewusst ein vom System under Test verschiedenes Modell gewählt, um
einen Selbstbewertungs\-Bias zu vermeiden, bei dem ein Modell anteilig seine eigenen
Ausgaben bewertet \cite{QUELLE_LLM_JUDGE}. Simulator und Judge laufen auf getrennten Endpunkten.
% =========================================================
% ABBILDUNG · Ablauf der Dialog-Simulation
% =========================================================
\begin{figure}[htbp]
  \centering
  \begin{tikzpicture}[
    node distance=10mm and 14mm,
    box/.style={draw, rounded corners, align=center, minimum height=10mm,
                inner sep=4pt, font=\small},
    lbl/.style={font=\scriptsize},
    >={Stealth[]}
  ]
    \node[box] (golden) {Golden\\(Szenario + feste\\Initialfrage)};
    \node[box, right=of golden] (sim) {ConversationSimulator\\\emph{Studenten-Simulator}};
    \node[box, right=of sim] (rag) {RAG-Tutor\\(System under Test)};
    \node[box, below=14mm of sim] (tc) {ConversationalTestCase\\(Multi-Turn-Dialog)};
    \node[box, right=of tc] (judge) {Judge\\(Bewertung)};
 
    \draw[->] (golden) -- (sim);
    \draw[->] ([yshift=2mm]sim.east) -- node[lbl, above]{Frage} ([yshift=2mm]rag.west);
    \draw[->] ([yshift=-2mm]rag.west) -- node[lbl, below]{Antwort} ([yshift=-2mm]sim.east);
    \draw[->] (sim) -- (tc);
    \draw[->] (tc) -- (judge);
  \end{tikzpicture}
  \caption{Ablauf der Dialog-Simulation: Aus einem Golden erzeugt der
  ConversationSimulator (in der Rolle der studierenden Person) fortlaufend
  Fragen an den RAG-Tutor; der resultierende Multi-Turn-Dialog wird anschlie{\ss}end
  vom Judge bewertet.}
  \label{fig:simulationsablauf}
\end{figure}
\subsection{Verwendete Methoden und Werkzeuge}
\label{subsec:methoden}

Die technische Umsetzung erfolgt in Python unter Verwendung des
Evaluations\-Frameworks DeepEval. Die Multi\-Turn\-Gespräche
werden mit dem ConversationSimulator des Frameworks erzeugt. Ausgehend von der
festen Initialfrage generiert der Simulator über eine festgelegte Anzahl von Runden
fortlaufend Studierenden\-Nachrichten, die über eine Callback\-Schnittstelle an die
RAG\-Pipeline gesendet werden. Die Antwort bildet jeweils den nächsten
Assistant\-Turn. Die Zuordnung von Antworten zu Metriken erfolgt anschließend über
ConversationalTestCase\-Objekte des Frameworks.

Für die kontextbezogenen Metriken wird der tatsächlich abgerufene Retrieval\-Kontext
benötigt. Dieser wird über den Verbose\-Modus der RAG\-Pipeline gewonnen, der bei
deaktiviertem Streaming eine strukturierte Antwort liefert, aus der die finalen
Retrieval\-Chunks samt Text extrahiert und den jeweiligen Antworten zugeordnet werden.
Zwei Aufbereitungsschritte werden dabei transparent gemacht: Zum einen wird der
Kontext für die Bewertung auf eine handhabbare Länge begrenzt, um ein Überschreiten des
Kontextfensters des Judge zu vermeiden; zum anderen wird die mathematische Notation
(LaTeX) für die Bewertung normalisiert, da sie andernfalls die strukturierte Ausgabe
des Judge stört. Beide Eingriffe betreffen ausschließlich die dem Judge übergebene
Kopie; die gespeicherten Gespräche bleiben unverändert.

\subsection{Bewertungskriterien}
\label{subsec:kriterien}

Die Bewertung ist entlang der beiden Forschungsfragen unterschiedlich ausgestaltet.

\subsubsection{Sokratische Qualität (RQ1)}

Sokratisches Tutoring ist ein didaktik\-spezifisches Konstrukt, für das im verwendeten
Framework keine vorgefertigten Metriken existieren. Es wird daher über eigene, mittels
G\-Eval operationalisierte Kriterien bewertet, die jeweils über
explizite Bewertungsschritte und eine dreistufige Bewertungsskala (Rubric) verankert
sind. Verwendet werden die Kriterien \emph{Lösung nicht verraten}, \emph{Sokratische
Rückfragen}, \emph{Schrittweise Lernprogression}, \emph{Niveau\-Anpassung},
\emph{Fachliche Korrektheit}, \emph{Verständnisförderung/Transfer} sowie
\emph{Respektvolle Kommunikation}. Ergänzend wird die im Framework enthaltene Metrik
zur Rollentreue (\emph{RoleAdherence}) als Gegenprüfung herangezogen, da sie die Treue
zur im Testfall hinterlegten sokratischen Rollendefinition misst und damit als einzige
native Metrik das Konstrukt indirekt berührt.

Bewusst nicht verwendet werden Metriken, die dem sokratischen Ziel widersprechen oder es
verfehlen. Dies betrifft insbesondere eine Metrik zur direkten Zielerreichung, die die
unmittelbare Beantwortung der Fachfrage bewertet und damit dem Prinzip, die Lösung
gerade nicht vorwegzunehmen, entgegensteht, sowie eine generische Turn\-Relevanz\-Metrik,
die legitimes sokratisches Umlenken als weniger relevant einstufen kann.

Zur Beantwortung von RQ1 werden vier Prompt-Varianten verglichen, die sich im Grad der didaktischen Steuerung unterscheiden. Der ausführliche sokratische System-Prompt (\texttt{system\_prompt}) spezifiziert die Gesprächsregeln explizit nach dem Framework der Universität. Der minimalistisch-sokratische Prompt (\texttt{minimaler\_sokrat}) enthält nur die Kernanweisung, ohne die Regeln auszuformulieren. Eine Variante ohne steuernden Prompt (\texttt{no\_Prompt}) dient als Referenz ohne sokratische Vorgabe. Die Standard-Variante (\texttt{stock\_prompt}) verwendet einen generischen Tutor-Prompt ohne sokratische Ausrichtung. Die Varianten spannen damit ein Kontinuum von detaillierter über minimale bis hin zu keiner sokratischen Instruktion auf, wodurch sich der Einfluss des Promptings auf die sokratische Qualität untersuchen lässt. Der vollständige Wortlaut der Prompts ist im Anhang dokumentiert.


% =========================================================
% TABELLE 1 · Uebersicht der RQ1-Bewertungskriterien + Fundierung
% =========================================================
\begin{table}[htbp]
    \centering
    \caption{Bewertungskriterien der sokratischen Qualität (RQ1) mit ihrer
    literaturgestützten Fundierung. Sechs Kriterien werden als eigene
    ConversationalGEval-Metriken operationalisiert; die Rollentreue (RoleAdherence)
    wird als native Metrik ergänzend als Gegenprüfung herangezogen.}
    \label{tab:rq1-kriterien}
    \small

    \begin{tabular}{p{3.6cm}p{6.4cm}p{3.4cm}}
        \toprule
        \textbf{Kriterium} &
        \textbf{Gemessenes Konstrukt} \\
        \midrule

        Lösung nicht verraten &
        Führt zur eigenen Erkenntnis, statt die Lösung zu nennen; lässt eigene Denkarbeit. \\

        \addlinespace

        Sokratische Rückfragen &
        Zielgerichtete, offene Fragen -- keine Fragen ,,um der Frage willen``.\\

        \addlinespace

        Schrittweise Lernprogression &
        Neues Wissen in kleinen Schritten auf Vorwissen aufbauend. \\

        \addlinespace

        Niveau-Anpassung &
        Diagnostiziert das Niveau und unterrichtet daran angepasst. \\

        \addlinespace

        Fachliche Korrektheit &
        Real korrekte Fachinhalte; keine erfundenen Formeln/Werte.\\

        \addlinespace

        Verständnisförderung / Transfer &
        Betont wichtige Lernpunkte; fördert Verständnis statt Reproduktion.\\

        \addlinespace

        Respektvolle Kommunikation &
        Ziel ist Lehren, nicht Bloßstellen; keine Beschämung. \\

        \midrule

        RoleAdherence (nativ) &
        Gegenprüfung der Treue zur sokratischen Rollendefinition.\\

        \bottomrule
    \end{tabular}
\end{table}

\subsubsection{Retrieval\-Qualität (RQ2)}

Für den Vergleich der RAG\-Konfigurationen werden die nativen Metriken
\emph{TurnContextualRelevancy} (Passung des abgerufenen Kontexts zur Frage) und
\emph{TurnFaithfulness} (Deckung der Antwort durch den abgerufenen Kontext) verwendet.
Beide beziehen den zuvor extrahierten Retrieval\-Kontext ein und eignen sich damit zur
Unterscheidung von Retrieval\-Konfigurationen.
% Benötigte Pakete (Praeambel):
% \usepackage{booktabs}
% \usepackage{array}

% =========================================================
% TABELLE 1 · RQ2-Bewertungskriterien (Retrieval-Qualitaet)
% =========================================================
\begin{table}[htbp]
  \centering
  \caption{Bewertungskriterien der Retrieval-Qualit\"at (RQ2). Die beiden Metriken sind
  native DeepEval-Metriken und beziehen den tats\"achlich abgerufenen
  \texttt{retrieval\_context} ein; die Abdeckung und die effektive Retrieval-G\"ute
  werden daraus abgeleitet.}
  \label{tab:rq2-kriterien}
  \small
  \begin{tabular}{p{4.0cm}p{2.0cm}p{7.0cm}}
    \toprule
    \textbf{Kennzahl} & \textbf{Herkunft} & \textbf{Gemessenes Konstrukt} \\
    \midrule
    TurnContextualRelevancy & nativ & Passung des abgerufenen Kontexts zur jeweiligen Frage. \\
    \addlinespace
    TurnFaithfulness        & nativ & Deckung der Antwort durch den abgerufenen Kontext (Grounding). \\
    \addlinespace
    Retrieval-Abdeckung     & abgeleitet & Anteil der Antworten mit tats\"achlich abgerufenem Kontext (aus den Konversationsdaten). \\
    \addlinespace
    Effektive Retrieval-G\"ute & abgeleitet & Produkt aus Abdeckung und Kontext-Relevanz; ber\"ucksichtigt, wie oft \emph{und} wie gut Kontext abgerufen wird. \\
    \bottomrule
  \end{tabular}
\end{table}

% =========================================================
% TABELLE 2 · Untersuchte RAG-Konfigurationen (Testmatrix RQ2)
% =========================================================
\begin{table}[htbp]
  \centering
  \caption{Untersuchte RAG-Konfigurationen (RQ2). Variiert werden die einbezogenen
  Collections, der Retrieval-Modus und der Einsatz des Cross-Encoder-Rerankings; eine
  reine LLM-Konfiguration ohne Retrieval dient als Referenz.}
  \label{tab:rq2-konfigurationen}
  \small
  \begin{tabular}{lp{4.2cm}p{2.6cm}c}
    \toprule
    \textbf{Konfiguration} & \textbf{Collections} & \textbf{Retrieval} & \textbf{Rerank} \\
    \midrule
    llm\_only                 & --                       & --            & Nein \\
    labor\_dense              & Labor                    & dense         & Nein \\
    vorlesung\_dense          & Vorlesung                & dense         & Nein \\
    beide\_dense              & Labor, Vorlesung         & dense         & Nein \\
    beide\_sparse             & Labor, Vorlesung         & sparse        & Nein \\
    beide\_hybrid             & Labor, Vorlesung         & dense, sparse & Nein \\
    beide\_hybrid\_rerank     & Labor, Vorlesung         & dense, sparse & Ja \\
    labor\_hybrid\_rerank     & Labor                    & dense, sparse & Ja \\
    vorlesung\_hybrid\_rerank & Vorlesung                & dense, sparse & Ja \\
    \bottomrule
  \end{tabular}
\end{table}
\subsubsection{Aggregation und Gate\-Logik}

Zur Beantwortung von RQ1 wird aus den G\-Eval\-Kriterien je Gespräch ein gewichteter
Gesamtscore gebildet [PLATZHALTER: konkrete Gewichte angeben und begründen]. Da eine
Lösungspreisgabe für einen sokratischen Tutor kein untergeordneter Mangel, sondern eine
Verletzung der Kernanforderung darstellt, werden die Kriterien \emph{Lösung nicht
verraten} und \emph{Fachliche Korrektheit} zusätzlich als primäre, nicht durch andere
Kriterien kompensierbare Anforderungen behandelt. Statt eines einfachen Mittelwerts,
der eine solche Verletzung durch gute Werte in anderen Kriterien kaschieren würde, wird
daher die Einhaltungsrate dieser Kernanforderungen als eigenständige Kennzahl berichtet.
Für RQ2 werden die Konfigurationen anhand der Mittelwerte der Retrieval\-Metriken sowie
der Retrieval\-Abdeckung, also des Anteils der Antworten mit tatsächlich abgerufenem
Kontext, verglichen.

\subsection{Durchführung}
\label{subsec:durchfuehrung}

Die Durchführung erfolgt zweiphasig und je Prompt\-Version bzw. Konfiguration getrennt:
Zunächst werden alle Gespräche simuliert, anschließend werden sie bewertet. Diese
Trennung erlaubt es, die netz\- und laufzeitintensive Simulation gegen den
RAG\-Endpunkt einmalig durchzuführen und die Judge\-Bewertung unabhängig davon zu
wiederholen. Zur Sicherstellung der Reproduzierbarkeit werden die simulierten Gespräche
je Konversation atomar zwischengespeichert; über eine deterministische Konversations\-ID
wird ein unterbrochener Lauf beim Neustart fortgesetzt, ohne bereits erzeugte Gespräche
neu zu berechnen. Die Bewertungsergebnisse werden fortlaufend je Metrik protokolliert
und ebenfalls fortsetzbar geschrieben.

Für RQ1 werden [PLATZHALTER: Anzahl] Prompt\-Varianten bei einer festen
RAG\-Konfiguration ausgewertet. Für RQ2 wird ein fester sokratischer Prompt mit einer
Matrix aus neun RAG\-Konfigurationen kombiniert, die sich in den einbezogenen
Collections, im Retrieval\-Modus (dense, sparse, hybrid) und im Einsatz des
Cross\-Encoder\-Rerankings unterscheiden und zusätzlich eine reine LLM\-Konfiguration
ohne Retrieval als Referenz umfassen. Zur Steuerung der Endpunkt\-Auslastung wird ein
gemeinsamer Rate\-Limiter eingesetzt; transiente Fehler (etwa Rate\-Limits oder
Timeouts) werden mit ansteigender Wartezeit wiederholt, deterministische Fehler
hingegen unmittelbar protokolliert.

Die Auswertung erfolgt anschließend automatisiert: Aus den protokollierten Rohdaten
werden je Auswertungslinie die beschriebenen Kennzahlen berechnet und zur weiteren
Analyse [PLATZHALTER: verwendete Auswertungs\-/Visualisierungswerkzeuge, z.\,B. pandas]
aufbereitet.

\subsection{Gütekriterien und methodische Einschränkungen}
\label{subsec:limitationen}

Sämtliche Bewertungen sind Urteile eines LLM\-Judge und damit modellabhängig; Ergebnisse
verschiedener Judge\-Modelle sind nicht unmittelbar vergleichbar. Die selbst definierten
G\-Eval\-Kriterien sind projektspezifisch und wurden [PLATZHALTER: Angabe, ob und wie
eine Validierung, z.\,B. eine manuelle Nachbewertung einer Stichprobe, durchgeführt
wurde] abgesichert. Simulator und RAG\-Tutor beruhen auf demselben Basismodell, was
durch die auf dem verwendeten Endpunkt verfügbaren Modelle bedingt ist; die methodisch
zentrale Trennung zwischen Judge und System under Test bleibt hingegen gewahrt. Da ein
Teil der Antworten ohne abgerufenen Kontext erzeugt wird, sind die kontextbezogenen
Metriken nur eingeschränkt aussagekräftig und werden entsprechend gesondert behandelt.
Aufgrund der begrenzten Anzahl an Szenarien je Konfiguration sind die Ergebnisse als
explorativ zu verstehen.